\chapter{Recommendation Systems}
\label{chap:recommendation}
\label{chap:recommendation_systems}

\begin{tldr}
Recommendation is deep learning's largest production deployment, and this book's full treatment of it now lives in Volume III: models and systems in [SR 6], evaluation and experimentation in [SR 7], serving and freshness in [SR 8]. This bridge chapter keeps what a Volume I reader needs: the one-arc view of recommendation as embeddings + metric learning + ranking, the retrieval--ranking funnel, the ten breadth facts every ML interview assumes, and two condensed high-frequency interview questions. If your interview loop is ads, ranking, or recommendations, study Volume III; if recommendation is one round among many, this chapter is your coverage.
\end{tldr}

\section{Recommendation as Deep Learning in One Arc}
\label{sec:recsys_bridge_arc}

Strip away the domain vocabulary and a recommender is a composition of techniques this volume has already built:

\begin{enumerate}
    \item \textbf{Embeddings} (Chapter~\ref{chap:embeddings}): users and items become vectors. In production the embedding tables---not the network---hold nearly all parameters: a DLRM-class ranking model reaches trillions of parameters, over 99\% of them embedding rows. The model is memory-bound by lookups, not compute-bound by matmuls.
    \item \textbf{Metric learning} (Chapter~\ref{chap:similarity_architectures}): retrieval trains a two-tower model with a sampled softmax / InfoNCE objective so that dot-product proximity means ``this user will engage with this item.'' The negatives come from the batch; a logQ popularity correction unbiases them.
    \item \textbf{Attention} (Chapter~\ref{chap:attention_transformers}): ranking models attend over the user's behavior sequence conditioned on the candidate item (DIN-style), and the 2024--26 generative turn (Meta's HSTU) models the raw action stream autoregressively, exactly as a decoder-only LM models tokens.
    \item \textbf{Ranking losses and calibration} (Chapter~\ref{chap:loss_functions}): pointwise cross-entropy on clicks, pairwise BPR, listwise softmax---and in ads, the predicted probability itself is a product requirement, because price is a function of pCTR.
    \item \textbf{Serving under a latency budget} (Chapter~\ref{chap:inference_optimization}): every architecture choice is negotiated against tens of milliseconds of end-to-end budget, which is what forces the funnel shape below.
\end{enumerate}

\section{The Retrieval--Ranking Funnel}
\label{sec:recsys_bridge_funnel}

No single model can score millions of items per request under a $\sim$100ms budget, so every production system is a funnel: cheap models over everything, expensive models over survivors.

\begin{figure}[H]
\centering
\begin{tikzpicture}[
    box/.style={rectangle, draw, minimum width=2.4cm, minimum height=0.8cm, rounded corners, font=\small, align=center}
]
    \node[box, fill=lightgray] (corpus) {Corpus\\{\footnotesize $10^6$--$10^9$ items}};
    \node[box, fill=lightblue, right=0.7cm of corpus] (retrieval) {Retrieval\\{\footnotesize $\rightarrow$ 1000s}};
    \node[box, fill=lightgreen, right=0.7cm of retrieval] (ranking) {Ranking\\{\footnotesize $\rightarrow$ 100s}};
    \node[box, fill=lightyellow, right=0.7cm of ranking] (rerank) {Re-rank\\{\footnotesize $\rightarrow$ 10s}};
    \draw[->] (corpus) -- (retrieval);
    \draw[->] (retrieval) -- (ranking);
    \draw[->] (ranking) -- (rerank);
    \node[below=0.2cm of retrieval, font=\footnotesize, text=deepblue, align=center] {two-tower\\+ ANN};
    \node[below=0.2cm of ranking, font=\footnotesize, text=deepblue, align=center] {DLRM / DCN /\\DIN};
    \node[below=0.2cm of rerank, font=\footnotesize, text=deepblue, align=center] {diversity,\\business rules};
\end{tikzpicture}
\caption{The retrieval--ranking funnel: per-item compute grows by orders of magnitude at each stage as the candidate set shrinks.}
\label{fig:recsys_bridge_funnel}
\end{figure}

\textbf{Retrieval} must be sub-linear in corpus size, so it uses the decomposed two-tower architecture: item embeddings pre-computed into an approximate nearest neighbor (ANN) index, only the user tower running online. \textbf{Ranking} can afford joint user-item computation---cross features, attention over behavior history, multi-task heads (CTR, CVR, engagement)---because it sees only hundreds of candidates. \textbf{Re-ranking} applies what the scoring model cannot see: list-level diversity, freshness boosts, exploration slots, policy constraints. The stages fail differently, too: retrieval failures are silent (a never-retrieved item cannot be ranked), while ranking failures are visible in metrics.

\section{Ten Facts to Know Cold}
\label{sec:recsys_bridge_facts}

\begin{enumerate}
    \item \textbf{Two-tower + ANN is the retrieval workhorse.} Decomposition is what buys corpus-scale serving; the price is expressiveness (dot product only), paid back at the ranking stage. Trained with in-batch sampled softmax; the logQ correction ($\text{logit} - \log Q(\text{item})$) fixes the popularity bias of in-batch negatives.
    \item \textbf{Interaction models rank; they cannot retrieve.} DLRM/DCN/DIN-class models need the user-item pair jointly and so can only score small candidate sets. ``Which model where'' is a latency-budget question, not a taste question.
    \item \textbf{Embedding tables dominate.} Trillion-parameter recommenders are trillions of \textit{embedding} parameters; the MLPs are megabytes. Sharding, mixed dimensions, and feature hashing are the systems levers.
    \item \textbf{Implicit feedback is abundant and biased.} Clicks are dense but confounded---position bias above all. Training on logged data without propensity correction teaches the model to imitate the old policy, which is why offline gains routinely vanish online.
    \item \textbf{Know NDCG@K, Recall@K, MRR} (Chapter~\ref{chap:evaluation_metrics}), and know that offline ranking metrics are weakly coupled to online business metrics: AUC can rise while revenue falls.
    \item \textbf{Calibration is revenue-critical in ads.} Rank score is $\text{pCTR} \times \text{bid}$ and pricing depends on the probability's absolute value, so a well-ranked but 2$\times$ over-confident model silently misprices the auction. AUC and calibration are orthogonal properties.
    \item \textbf{Cold start is a data problem.} New users: onboarding signals, context features, fast-adapting bandits. New items: content features and deliberate exploration. LLM-generated enrichment and content-derived semantic IDs are the modern cold-start on-ramp.
    \item \textbf{Exploration is mandatory.} The system trains on outcomes of its own choices; without an exploration budget (epsilon-greedy, UCB, Thompson sampling) and logged propensities, tail and fresh items never accumulate signal and off-policy evaluation is impossible.
    \item \textbf{Feedback loops are the field's signature failure mode.} Popularity bias amplifies, filter bubbles form, and engagement proxies drift from satisfaction (Goodhart's law); mitigation is diversity constraints, long-horizon objectives, and holdout traffic.
    \item \textbf{The 2024--26 shift is generative.} Meta's HSTU reframes recommendation as autoregressive sequence modeling over trillions of user action events (compute-bound, LM-style scaling laws, +12.4\% topline in production), and RQ-VAE semantic IDs (TIGER, YouTube) give items a content-derived token vocabulary---the bridge between LLM machinery and billion-item catalogs.
\end{enumerate}

\begin{keyinsight}[Where the Full Treatment Lives]
This chapter is a bridge. The complete material---collaborative filtering and matrix factorization, Wide\&Deep/DCN/DLRM/DIN, multi-task ads models (MMoE, PLE, ESMM), sequential and generative recommenders (SASRec through HSTU), semantic IDs and generative retrieval, LLMs in recommendation, bandits, the ads pipeline and auction, calibration, and embedding-table infrastructure---is Volume III, Chapter [SR 6]. Evaluation and experimentation for ranking systems is [SR 7]; serving, index freshness, and caching is [SR 8]. Within this volume, the load-bearing prerequisites are Chapters~\ref{chap:embeddings}, \ref{chap:similarity_architectures}, and \ref{chap:attention_transformers}.
\end{keyinsight}

\section{Interview Questions}
\label{sec:recsys_interview}

Deliberately none. Both questions this bridge used to carry---cold start, and
two-tower versus interaction-based candidate generation---exist verbatim in
[SR 6], which is the canonical home. Duplicating them here is exactly the drift
the program's canonical-home contract exists to prevent: two copies age apart,
and a reader cannot tell which is authoritative. The full set---system design,
estimation and debugging on recommenders---is [SR 6]; ranking evaluation is
[SR 7]; serving and freshness are [SR 8].
